\section{Описание}
Требуется реализовать алгоритм поиска образца, который может содержать
«джокеры»~--- символы `?`, совпадающие с любым числом из заданного
алфавита $\{0,\dots,2^{32}-1\}$. Для решения использован алгоритм
Ахо~--~Корасик, адаптированный для работы с маской.

Основная идея:
\begin{enumerate}
  \item Исходный образец разбивается на максимальные фрагменты,
        не содержащие джокеров.
  \item Для набора полученных фрагментов строится автомат
        Ахо~--~Корасик.
  \item Входной текст обрабатывается посимвольно; при обнаружении
        конца очередного фрагмента фиксируется позиция его окончания.
  \item Для каждого потенциального начала полного вхождения
        подсчитываются «голоса»: если все фрагменты обнаружены
        на ожидаемых позициях относительно этого начала, число голосов
        становится равным количеству фрагментов.
  \item Позиции, набравшие нужное количество голосов и не выходящие
        за границу текста, объявляются вхождениями.
\end{enumerate}

Сложность алгоритма складывается из трёх составляющих:
\begin{itemize}
  \item $O(n)$ --- посимвольная обработка текста автоматом (метод
        \texttt{feedSymbol} вызывается ровно $n$ раз, каждый вызов
        выполняется за амортизированное $O(1)$);
  \item $O(m)$ --- построение бора и суффиксных ссылок по всем
        фрагментам образца, где $m$~--- суммарная длина фрагментов;
  \item $O(Z)$ --- сохранение всех найденных окончаний фрагментов
        (метод \texttt{saveMatch}) и финальный подсчёт голосов
        (метод \texttt{getResults}), где $Z$~--- общее количество
        вхождений всех фрагментов в текст.
\end{itemize}

Таким образом, общая сложность составляет $O(n + m + Z)$.

В среднем и лучшем случаях, когда фрагменты образца встречаются в
тексте редко, величина $Z$ мала (сравнима с количеством итоговых
вхождений полного образца), и сложность близка к линейной $O(n + m)$.

Однако в худшем случае $Z$ может достигать $O(n \cdot m)$. Это
происходит, например, когда образец состоит из коротких фрагментов
(скажем, одиночных символов), разделённых джокерами, а текст
полностью состоит из символов, совпадающих с этими фрагментами.
В такой ситуации \emph{каждый} символ текста порождает вхождение
\emph{каждого} фрагмента, и $Z \approx n \cdot k$, где $k$~---
количество фрагментов ($k \leq m$). Тогда $Z = O(n \cdot m)$, и
общая сложность деградирует до $O(n \cdot m)$, что сравнимо с
наивным алгоритмом.

\section{Исходный код}
Программа реализована в одном файле \texttt{main.cpp}. Ниже перечислены
основные структуры данных и функции.

\subsection{Структуры данных}
\begin{lstlisting}[language=C++]
struct Symbol {
    unsigned int value;
    bool isJoker = false;

    Symbol(unsigned int value, bool isJoker)
        : value(value), isJoker(isJoker) {}
};

struct Pattern {
    std::vector<Symbol> symbols;

    Pattern() : symbols() {}
    Pattern(const std::vector<Symbol>& p) : symbols(p) {}

    friend std::istream& operator>>(std::istream& is, Pattern& pattern);
    friend std::ostream& operator<<(std::ostream& os, Pattern& pattern);
};

struct Match {
    size_t line;
    size_t pos;

    Match(size_t line, size_t pos) : line(line), pos(pos) {}
};

class AhoCorasick {
    struct Node {
        std::unordered_map<unsigned int, Node*> next;
        std::vector<size_t> fragment_offsets;
        Node* fail = nullptr;
        Node* out = nullptr;

        bool isEnd() const { return !fragment_offsets.empty(); }
        ~Node() { for (auto& [_, child] : next) delete child; }
    };

    Node* root;
    const Node* cur;
    std::vector<size_t> row_match_starts;
    size_t cur_pos = 0;
    std::vector<size_t> line_starts;
    size_t fragment_number = 0;
    size_t pattern_length = 0;

    void buildLinks();
    void addFragment(const std::vector<unsigned int>& fragment,
                     size_t i_end);
    void saveMatch(size_t abs_end, const Node* node);
    Match getLineAndColumn(size_t start_pos);

public:
    AhoCorasick(const Pattern& pattern);
    ~AhoCorasick() { delete root; }

    void feedSymbol(unsigned int symbol);
    void feedNewline();
    std::vector<Match> getResults();
};
\end{lstlisting}

\subsection{Основные методы и их назначение}
\begin{longtable}{|p{7.5cm}|p{7.5cm}|}
\hline
\rowcolor{lightgray}
\multicolumn{2}{|c|}{\texttt{main.cpp}} \\
\hline
\texttt{AhoCorasick(const Pattern\& pattern)} &
Конструктор: разбивает образец на фрагменты без джокеров, строит бор,
добавляет фрагменты (метод \texttt{addFragment}), вызывает
\texttt{buildLinks}. \\
\hline
\texttt{void buildLinks()} &
Построение суффиксных ссылок (\texttt{fail}) и выходных ссылок
(\texttt{out}) стандартным обходом в ширину. \\
\hline
\texttt{void addFragment()} & \multirow{2}{=}{Добавление в бор одного фрагмента; в конечную вершину заносится смещение конца фрагмента в образце.} \\
\texttt{(const vector<unsigned int>\& fragment, size\_t i\_end)} & \\
\hline
\texttt{void feedSymbol(unsigned int symbol)} &
Переход по символу текста; при попадании в конечную вершину (или в
вершину, из которой есть выход) вызывает \texttt{saveMatch} для
регистрации окончаний фрагментов. \\
\hline
\texttt{void feedNewline()} &
Запоминает текущую абсолютную позицию как начало новой строки. \\
\hline
\texttt{std::vector<Match> getResults()} &
После обработки всего текста подсчитывает голоса за каждую возможную
стартовую позицию, выбирает те, где число голосов равно числу
фрагментов и вхождение не выходит за границу текста. Результат
пересчитывается в пары (строка, позиция в строке). \\
\hline
\texttt{int main()} &
Считывает образец через оператор \texttt{>>}, создаёт автомат, затем
побайтово читает входной поток, распознаёт числа и переводы строк,
вызывает \texttt{feedSymbol} и \texttt{feedNewline}. В конце печатает
номера строк и позиций найденных вхождений. \\
\hline
\end{longtable}

\subsection{Фрагмент главного цикла чтения}
Для обработки чисел, разделённых пробельными символами, используется
конечный автомат с двумя состояниями:
\begin{lstlisting}[language=C++]
enum Status { waiting, reading_number };
Status status = waiting;
unsigned int value = 0;

while ((c = getchar()) && c != EOF) {
    switch (status) {
    case waiting:
        if (std::isdigit(c)) {
            value = c - '0';
            status = reading_number;
        }
        break;
    case reading_number:
        if (std::isdigit(c)) {
            value = 10 * value + c - '0';
        } else {
            ac.feedSymbol(value);
            status = waiting;
        }
        break;
    }
    if (c == '\n') ac.feedNewline();
}
\end{lstlisting}
\pagebreak

\section{Консоль}
\begin{alltt}
root@0d612ef3b342:/workspaces/discran-labs/laba-4# cat test.txt
1 ? 02
0001
1
02
2
root@0d612ef3b342:/workspaces/discran-labs/laba-4# make
g++ -std=c++20 main.cpp -o main.out
root@0d612ef3b342:/workspaces/discran-labs/laba-4# ./main.out < test.txt
1, 1
2, 1
\end{alltt}
\pagebreak

